\makeatletter
\def\input@path{{../00/plantilla/fisica/}}
\makeatother
%% ============================================================================
%%  LAFROX-Formulario-T-11.tex
%%
%%  Formulario T-11: Especificaciones técnicas ofertadas. Archivo propio, separado
%%  del Subdocumento 4 (aclaraciones, sección 1: un formulario por archivo).
%%  Se compila con compilar.ps1 en la raíz del proyecto.
%%  El cuerpo está en formularios/T11/15_anexo_t11.tex.
%%
%%  Columnas oficiales (Bases Administrativas, Formulario T-11): Componente |
%%  Producto / servicio ofertado | Ubicación / Lugar | Cantidad | Justificación.
%%  Sin precios (Art. 50.2).
%% ============================================================================
\documentclass[carta,firma]{lafrox}

\makeatletter
\@removefromreset{table}{chapter}
\makeatother
\renewcommand{\thetable}{\arabic{table}}

% Tablas sin palabras cortadas (aclaraciones, sección 5).
\let\LFXcuerpotablaorig\LFXcuerpotabla
\renewcommand{\LFXcuerpotabla}{%
  \LFXcuerpotablaorig\hyphenpenalty=10000\exhyphenpenalty=10000\relax}

\datoslafrox{
  documento  = Propuesta Técnica,
  subtitulo  = Formulario T-11 --- Especificaciones técnicas ofertadas,
  alcance    = {Subdocumento 4 --- Arquitectura lógica y física de la solución},
  formulario = {Formulario T-11},
  fecha      = 05 de octubre de 2026,
  version    = {},
  nomenclatura = {LAFROX-Formulario-T-11.pdf},
}
\usepackage{etoolbox}
\renewcommand{\LFXestadodoc}{Entrega 2}
\patchcmd{\portadalafrox}{\LFXcampoportada{Versión}{\LFXversion}}{\LFXcampoportada{Documento}{Formulario T-11}}{}{}
\patchcmd{\LFXpieizquierdo}{\textperiodcentered\ v\LFXversion}{}{}{}

\begin{document}
\portadalafrox

% El cuerpo del formulario vive en formularios/T11/15_anexo_t11.tex y se compila solo
% en este archivo propio (aclaraciones, sección 1). Se edita solo ahí.
% !TeX root = ../../main.tex
% ------------------------------------------------------------------
%  Formulario T-11: especificaciones técnicas ofertadas.
% ------------------------------------------------------------------

\chapter*{Formulario T-11: Especificaciones técnicas ofertadas}
\markboth{Formulario T-11}{Formulario T-11}
\phantomsection\addcontentsline{toc}{chapter}{Formulario T-11: Especificaciones técnicas ofertadas}

Este formulario detalla los componentes de infraestructura, plataforma, licenciamiento y hardware especificados para la solución, con su ubicación, su cantidad y su justificación, conforme al Formulario T-11 de las Bases Administrativas. Acompaña al Subdocumento 4 (Arquitectura lógica y física de la solución): su resumen y análisis están en el apartado 4.2.1 de ese subdocumento, y el emplazamiento de cada componente se justifica en su apartado 4.2.2. Cuando un elemento materializa un componente del catálogo de emplazamiento, la columna Componente lleva el mismo código y nombre del apartado 4.2.2 (por ejemplo, D-01 Firewall y Customer Gateway); los elementos de soporte que no son componentes de arquitectura, como la infraestructura de la sala técnica o el cableado, se identifican solo por su nombre. El hardware de terreno ---terminales, lectores, impresoras portátiles y sensores--- lo adquiere el CLIENTE, y el adjudicatario lo especifica (Bases Técnicas del caso, cap. 11, p. 20). La infraestructura on-premise la provee el adjudicatario dentro del contrato (Bases Administrativas, art. 14.2, p. 10). Todo el equipamiento lo instala, integra y mantiene el adjudicatario.

La Tabla~\ref{tab:t11} reúne todos los elementos ofertados, agrupados por familia: infraestructura de cómputo y almacenamiento, red y seguridad, dispositivos de terreno y de operación, sala técnica de Talca, gabinetes de borde, plataforma de nube y software de base. Las cantidades de dispositivos de terreno y de componentes críticos incluyen una reserva del 10 \% del parque de cada tipo, redondeada hacia arriba, conforme a la tabla de repuestos del numeral 8.4 (Bases Técnicas Transversales, cap. 8, p. 19); los supuestos que fijan cada cantidad (S-28 y S-30 a S-41) se registran en el Subdocumento 3, y el cálculo de cada cantidad está en el Anexo 4-W.

\begin{tablalafrox}{Formulario T-11: especificaciones técnicas ofertadas}{tab:t11}%
  {F{0.17} Y F{0.16} G{0.13} Y}%
  {\cab{Componente} & \cab{Producto / servicio ofertado} & \cab{Ubicación / Lugar} & \cab{Cantidad} & \cab{Justificación}}
  \multicolumn{5}{@{}l}{\intersemibold Infraestructura de cómputo y almacenamiento} \\
Servidor del clúster & HPE ProLiant DL345 Gen11 o Dell PowerEdge R7615; AMD EPYC 9124 de 16 núcleos y 32 hilos, 64 GB DDR5 ECC (4 × 16 GB), dos M.2 de 480 GB en RAID 1 para arranque, dos NVMe de 960 GB para Ceph sin RAID en modo HBA, dos puertos de 10 GbE y dos fuentes redundantes de 800 W en circuitos A/B & CD Talca, rack R01 & 3 & Los tres nodos conservan quórum y capacidad a 3× tras la pérdida de uno; Ceph utiliza seis OSD, tres réplicas y mínimo dos réplicas (RT-03.14; Bases Técnicas Transversales, cap. 3, p. 9; RT-08.04; Bases Técnicas Transversales, cap. 8, p. 18). \\
  D-05 --- Respaldo local & Synology RackStation RS2423RP+ o equivalente, fuente redundante, potencia de placa de 400 W, cuatro discos de 4 TB en RAID 6, aproximadamente 8 TB útiles, cifrado e instantáneas inmutables WORM & CD Talca, rack R01 & 1 & Constituye la copia local de recuperación rápida del esquema 3-2-1-1-0 y complementa la copia inmutable en nube. \\
  Medio de respaldo físico rotativo & Medio transportable cifrado, con rotación semanal y servicio de custodia externa & CD Talca y custodia externa & 2 medios: uno en sitio y otro en custodia; 1 servicio de custodia & Mantiene una copia física fuera del sitio y permite intercambiar semanalmente el medio bajo cadena de custodia (RT-06.26; Bases Técnicas Transversales, cap. 6, p. 16). \\
  Traslado del servidor del ERP & Desmontaje, traslado y montaje del servidor existente del ERP de 2017 del CLIENTE, sin cambios de software, tras un respaldo completo verificado; cota eléctrica de 1.600 W hasta medirlo & Desde la sala actual de 25 m\textsuperscript{2} al CD Talca, rack R01 (U25--U26) & 1 servicio & La sala actual no cumple el Capítulo 6 y se libera (Bases Técnicas del caso, cap. 17, p. 33); el ERP que emite las guías queda con UPS N+1, generador y climatización redundante (apartado 4.3.1.4). \\
  Conmutador de transferencia de rack & Conmutador automático de transferencia de 1U entre los circuitos A y B & CD Talca, rack R01 (U27) & 1 & Da doble alimentación al servidor del ERP si este tiene una sola fuente (RT-08.04; Bases Técnicas Transversales, cap. 8, p. 18); su consumo queda en el margen del 20 \% de la tabla de carga eléctrica del apartado 4.3.1. \\
  Servidores de Concepción (activo y en espera) & HPE ProLiant DL20 Gen11 con Xeon E-2400 de 8 núcleos y 16 hilos, 32 GB RAM, cuatro SSD de 1,92 TB en RAID 10 y dos fuentes redundantes de 500 W en circuitos distintos; consumo de placa de 1.000 W con ambas fuentes por servidor & CD Concepción, gabinete de borde & 2 & Sostienen el WMS autónomo de su propia bodega: el activo ejecuta VM-C01 a VM-C04 y el otro mantiene sus copias en espera, con la base replicada en forma sincrónica, y toma el control si el activo falla (RT-03.14; Bases Técnicas Transversales, cap. 3, p. 9); la capacidad se fundamenta en la tabla de VMs del Anexo 4-W, su consumo entra en el cálculo de la UPS del gabinete (RT-08.01; Bases Técnicas Transversales, cap. 8, p. 18) y la doble fuente cumple RT-08.04 (Bases Técnicas Transversales, cap. 8, p. 18). \\
  E-01 --- WMS de cross-docking (activo y en espera) & Mini-PC industrial Advantech ARK-2250 o equivalente con dos bahías de almacenamiento, sin ventilador y con operación hasta 60\,\textdegree{}C, Intel Core i5 o i7, 16 GB y dos SSD industriales de 512 GB en RAID 1 por software; consumo máximo de 45 W; dos fuentes DIN de 24 VDC en circuitos distintos y módulo de redundancia & Curicó, Chillán y Los Ángeles, 2 por plataforma; reserva en Talca & 6 + 1 de reserva = 7 & Opera el perfil WMS local de cada plataforma: el activo atiende y el otro mantiene la base replicada en forma sincrónica y toma el control si el activo falla (RT-03.14; Bases Técnicas Transversales, cap. 3, p. 9). Cada equipo tiene una interfaz a cada switch; una fuente toma la UPS del gabinete y otra la red, y la unidad de reserva se conserva en Talca. \\
  \multicolumn{5}{@{}l}{\intersemibold Red y seguridad} \\
  D-01 --- Firewall y Customer Gateway (Talca) & Par activo/pasivo FortiGate 100F o equivalente, con sincronización de sesiones, doble fuente y potencia de placa de 150 W por equipo & CD Talca, rack R02; reserva en Talca & 2 + 1 de reserva = 3 & El par conserva el perímetro y los túneles; cada equipo usa fuentes en circuitos distintos (RT-08.03 y RT-08.04; Bases Técnicas Transversales, cap. 8, p. 18). La reserva es intercambiable con el par de Concepción. \\
  D-02 --- Switching (núcleo de Talca) & Cisco Catalyst 9300-48P o equivalente en stack, con doble fuente y potencia de placa de 150 W por equipo & CD Talca, rack R02; reserva en Talca & 2 + 1 de reserva = 3 & Los dos equipos forman un único plano de distribución y cada uno tiene doble fuente (RT-08.03 y RT-08.04; Bases Técnicas Transversales, cap. 8, p. 18). La unidad de reserva cubre los switches de núcleo de Talca y de Concepción. \\
  D-02 --- Switching (gestión de Talca) & Switch de capa 2 con doble fuente y potencia de placa de 50 W para la VLAN de gestión & CD Talca, rack R02 & 1 & Aísla los puertos de administración de los servidores en una VLAN dedicada con MFA, y permite recuperar los equipos aunque la red de producción esté caída. \\
  Consola KVM & KVM sobre IP Raritan Dominion KX IV o equivalente, con doble fuente y potencia de placa de 50 W, cubierta por el margen eléctrico del 20 \% de la tabla de carga eléctrica del apartado 4.3.1 & CD Talca, rack R01 (U24) & 1 & Da acceso a la consola de los tres nodos y de la NAS mediante dos circuitos eléctricos. \\
  Distribución horizontal & Paneles y terminaciones de fibra OM4 y cobre Cat6A & CD Talca, racks R01 y R02 & 1 & Distribuye las conexiones terminadas en los racks. Las comunicaciones ingresan al edificio por dos puntos separados y siguen ductos independientes hasta la sala (RT-06.32; Bases Técnicas Transversales, cap. 6, p. 17). \\
  Red inalámbrica industrial & Puntos de acceso Wi-Fi 6E industriales en la VLAN de operaciones; los de la cámara de congelado, con rango de operación hasta $-$40\,\textdegree{}C & Talca 39, Concepción 20 y cross-docking 6 (2 por plataforma) & 65 + 7 de reserva = 72 & Estimación preliminar por superficie, con mayor densidad dentro de las cámaras (S-36); la cantidad definitiva la fija el estudio de cobertura. La red se segmenta por tipo de dispositivo, autentica con certificado o credencial de empresa y verifica cobertura dentro de las cámaras (RT-03.23; Bases Técnicas Transversales, cap. 3, p. 10). \\
  D-03 --- Enlace de fibra & Fibra óptica simétrica: Talca de 20 a 50 Mbps y Concepción de 10 a 20 Mbps & CD Talca y CD Concepción & 2 & Dimensionado para la replicación, el broker, la telemetría, la observabilidad y el incremental diario; admite el respaldo completo dominical de 15,12 GB a 20 Mbps en Talca y 7,76 GB a 10 Mbps en Concepción. El drenaje de 24 h en 2 h requiere 1,87 Mbps en Talca y 1,44 Mbps en Concepción; los peores casos son 5,16 y 2,14 Mbps, respectivamente, cubiertos por los rangos contratados de 20 a 50 y 10 a 20 Mbps (RT-03.20; Bases Técnicas Transversales, cap. 3, p. 9). \\
  D-04 --- Enlace LTE & LTE empresarial con SIM dedicada: Talca de 5 a 10 Mbps, Concepción de 3 a 5 Mbps y cross-docking de 2 a 5 Mbps con dos proveedores & Talca 1, Concepción 1 y 2 por cross-docking & 8 planes & Aporta caminos y proveedores distintos (RT-03.17; Bases Técnicas Transversales, cap. 3, p. 9), con conmutación en menos de 30 s. En los CD respalda a la fibra y antecede a Starlink, que queda en espera caliente; en los cross-docking, dos proveedores LTE respaldan el camino principal Starlink. \\
  Plan de datos móviles de terreno & SIM de datos empresarial, administrada por el MDM & Terminales EC55 de preventa y TC58e de reparto & 77 + 121 = 198 SIM, incluida la reserva y el crecimiento del año 3 & Sincroniza las capturas de preventa y reparto por red celular hacia CloudFront y recibe las actualizaciones de aplicación y sistema operativo de la preventa, que no pasan por los enlaces de los centros. \\
  D-01 --- Firewall y Customer Gateway (túneles VPN) & AWS Site-to-Site VPN IPsec IKEv2 con BGP sobre Transit Gateway & Talca, Concepción y los tres cross-docking hacia sa-east-1; configuración preparada también hacia us-east-1 & 2 túneles por sitio y región: 10 hacia sa-east-1 y 10 preparados hacia us-east-1 & Cifra los flujos de replicación, mensajería y telemetría, con enrutamiento dinámico y MTU de túnel de 1.436 bytes (RT-03.21; Bases Técnicas Transversales, cap. 3, p. 9). Es el mismo conjunto de túneles contado en la fila D-01 de la parte en nube. \\
  D-01 --- Firewall y Customer Gateway (Concepción) & Par activo/pasivo FortiGate 100F o equivalente, con sincronización de sesiones, doble fuente y potencia de placa de 150 W por equipo & CD Concepción, gabinete de borde & 2 (reserva compartida con Talca) & El par protege la bodega autónoma y conmuta entre fibra, LTE y Starlink en espera caliente; sus fuentes ocupan circuitos distintos (RT-08.03 y RT-08.04; Bases Técnicas Transversales, cap. 8, p. 18). \\
  D-02 --- Switching (núcleo de Concepción) & Cisco Catalyst 9300-48P o equivalente en stack, con doble fuente y potencia de placa de 150 W por equipo & CD Concepción, gabinete de borde & 2 (reserva compartida con Talca) & El par mantiene la distribución de la bodega autónoma ante la falla de un switch; cada equipo utiliza doble fuente (RT-08.03 y RT-08.04; Bases Técnicas Transversales, cap. 8, p. 18). \\
  D-02 --- Switching (acceso de bodegas) & Cisco Catalyst C9200-24P o equivalente, PoE+, enlace de fibra y segunda fuente & Gabinetes de piso: Talca 4 y Concepción 2; reserva en Talca & 6 + 1 de reserva = 7 & Conecta los puntos de acceso y periféricos cercanos al núcleo por fibra; la segunda fuente y la UPS del gabinete protegen la continuidad eléctrica. \\
  D-02 --- Switching (cross-docking) & Dos switches industriales gestionables PoE+ de doble entrada por plataforma, Moxa EDS-P506E-4PoE o equivalente; dos fuentes DIN por switch en circuitos distintos & Curicó, Chillán y Los Ángeles; reserva en Talca & 6 + 1 de reserva = 7 & Cada uno de los dos mini-PC se conecta a ambos switches y los dos puntos de acceso se distribuyen entre ellos. Cada switch reúne en sus seis puertos un punto de acceso, los dos mini-PC, los firewalls y el enlace entre switches. El par cumple RT-08.03 y su alimentación cumple RT-08.04 (Bases Técnicas Transversales, cap. 8, p. 18). \\
  D-01 --- Firewall y Customer Gateway (cross-docking) & Par FortiGate Rugged 60F-3G4G o equivalente, con IPsec IKEv2, BGP, módem LTE de doble SIM y doble entrada DC por equipo & Curicó, Chillán y Los Ángeles, 2 por plataforma; reserva en Talca & 6 + 1 de reserva = 7 & Termina el camino satelital y dos caminos LTE, con alimentación desde circuitos distintos y conmutación en menos de 30 s (RT-03.17, RT-08.03 y RT-08.04). \\
  \multicolumn{5}{@{}l}{\intersemibold Dispositivos de terreno y de operación} \\
C-01 --- App de preventa (terminal) & Zebra EC55: Android, IP67, uso con guantes y lector 2D SE4100; accesorios: cargador y funda con correa de mano & Uno por preventista; 70 en el año 3 & 70 + 7 de reserva = 77 & Opera sin señal, dura la jornada con una carga, funciona con guantes y su pantalla se lee a pleno sol (RT-08.11; Bases Técnicas Transversales, cap. 8, p. 19); la reserva es el 10 \% del parque. \\
  C-02 --- App de reparto (terminal) & Zebra TC58e: 5G, lector SE55, batería de 7.000 mAh, GPS L1/L5, IP65/68 y operación de $-$20 a +50\,\textdegree{}C; accesorios: soporte y cargador de vehículo & Uno por camión: base actual de 42 propios y 54 de transportistas; 110 camiones proyectados al año 3 & 110 + 11 de reserva = 121 & Incluye el crecimiento del año 3; se asigna al camión y no al conductor, porque los conductores de terceros rotan sin aviso (S-30); un solo modelo para conductores propios y externos: mismo repuesto, misma capacitación y misma integración. \\
  Impresora de cabina & Zebra ZQ620 Plus: 3 pulgadas, ZPL, IP54 y Link-OS; accesorios: batería de reemplazo y soporte de vehículo; consumible: rollo térmico & Una por camión; 110 en el año 3 & 110 + 11 de reserva = 121 & Imprime la evidencia de entrega y los comprobantes en la cabina; se asigna al camión (S-30). \\
  Terminal de pago & PAX A920 Pro: Android 14, PCI PTS 7.x y EMV L1/L2; accesorio: base de carga & Uno por camión; 110 en el año 3 & 110 + 11 de reserva = 121 & Cobra con tarjeta en el local del cliente con certificación PCI PTS y EMV, integrado por Bluetooth con la aplicación de reparto; se asigna al camión (S-30). \\
  C-03 --- Terminales de bodega (congelado) & Zebra MC9400 Cold Storage Freezer: $-$30\,\textdegree{}C, lector SE58 de 100 pies, batería de 5.000 mAh, Wi-Fi 6E, protección IP65 e IP68 y resistencia a caídas de 3 m sobre concreto; accesorios: batería de repuesto por equipo y cargador de baterías & CD Talca, cuadrilla de congelado & 23 + 3 de reserva = 26 & Prepara pedidos en la cámara a $-$22\,\textdegree{}C con guantes gruesos, lee a distancia y recambia la batería en caliente (RT-08.11 y RT-08.12; Bases Técnicas Transversales, cap. 8, p. 19; RT-13.08; Bases Técnicas del caso, cap. 15, p. 27); lo usa solo la cuadrilla que entra al congelado, cuyo tamaño sigue a la participación del congelado en las líneas (S-34). \\
  C-03 --- Terminales de bodega (estándar) & Zebra MC9400 estándar: lector SE58, batería de 5.000 mAh y Wi-Fi 6E; accesorio: cuna de carga de cinco posiciones & Talca 113, Concepción 68 y cross-docking 6 (2 por plataforma); reserva: Talca 12, Concepción 7 y una en cada plataforma & 187 + 22 de reserva = 209 & Incluye el crecimiento del año 3; atiende la bodega seca, la cámara de refrigerado sobre 0\,\textdegree{}C y la desconsolidación de los cross-docking; se comparte entre turnos (RT-12.11; Bases Técnicas del caso, cap. 15, p. 27), por lo que lo dimensiona el turno con más personas a la vez: 120 preparadores en Talca, 60 en Concepción (S-39) y 2 por cross-docking (S-35), sin personal de carga adicional (S-33). \\
  C-04 --- Impresoras de andén & Zebra ZT411: 203 dpi, 14 ips, Ethernet y ZPL; consumibles: etiquetas de 4 × 6 pulgadas y cinta de transferencia térmica & Talca 4 y Concepción 2; reserva en Talca & 6 + 1 de reserva = 7 & Imprime la etiqueta SSCC GS1-128 al armar la unidad logística (RF-01.07), legible a lo largo de toda la cadena según el estándar sectorial del caso (RT-05.23; Bases Técnicas Transversales, cap. 5, p. 13; Bases Técnicas del caso, cap. 15, p. 26). \\
  C-05 --- Balanzas de recepción & Dibal BEV con plataforma ME e indicador DMI-610 inoxidable, RS-232 y Ethernet; accesorio: rampa de acceso para transpaleta & Talca 2 y Concepción 1; reserva en Talca & 3 + 1 de reserva = 4 & Captura el peso de recepción en línea hacia el WMS (RF-01.01), sin digitación, con calibración certificada. \\
  B-01 --- Sensores de temperatura & Módulo Ebyte ME31-XDXX0400: 4 canales PT100, Modbus RTU/TCP, de $-$40 a +85\,\textdegree{}C, con sondas de 3 hilos; dos fuentes de 24 VDC en circuitos distintos y caja IP65 & Cámaras de Talca (15 puntos) y Concepción (6 puntos) & 6 + 1 de reserva = 7 módulos; 21 + 3 de reserva = 24 sondas & Alimenta el registro continuo de la cadena de frío, con una exactitud de $\pm$1,1\,\textdegree{}C a $-$22\,\textdegree{}C que distingue una excursión real de la variación del instrumento; los puntos se distribuyen por superficie de cámara y junto a cada puerta, y la cámara de Concepción se estima en unos 450 m² (S-37). \\
  B-03 --- Termógrafos de camión & Onset InTemp CX450: Bluetooth, de $-$30 a +70\,\textdegree{}C con $\pm$0,5\,\textdegree{}C y calibración NIST en dos puntos; accesorio: soporte de montaje en la caja de frío; consumible: batería reemplazable & Camiones con equipo de frío: 18 propios y 10 de transportistas & 28 + 3 de reserva = 31 & Registra la temperatura del transporte de forma independiente y sin conexión durante toda la ruta; alerta la excursión térmica en el terminal del conductor y, al volver, descarga el registro completo por Bluetooth a ese terminal, que lo envía a la nube. Los 10 camiones refrigerados de los transportistas se instrumentan con el mismo equipo, con el acuerdo de cada transportista (S-28); la flota con frío no crece en el horizonte de cotización (S-38). \\
  Estación de trabajo & Ocho PC nuevas; equipos existentes de oficina con CrowdStrike Falcon, cifrado de disco, actualización automatizada y control de extraíbles & Talca 5 y Concepción 3; oficinas de Talca y Concepción & 8 nuevas; 184 existentes gestionadas (S-41) & Puestos de despacho, administración, planificación, calidad y TI, con dos monitores de 24 pulgadas, ergonomía NCh 2527 y certificación de eficiencia energética; la postura gestionada habilita Verified Access (RT-08.07--08.09; Bases Técnicas Transversales, cap. 8, p. 19). \\
  D-06 --- Enlace satelital & Starlink fijo, antena y terminal alimentado desde circuito respaldado; subida mínima supuesta de 2 Mbps para el cálculo & Talca 1, Concepción 1 y cross-docking 3 & 5 & En los CD es el tercer camino en espera caliente: encendido, con túnel IPsec establecido y BGP de menor preferencia; cursa tráfico solo si fallan fibra y LTE. En esa falla prioriza por calidad de servicio DMS/WAL de Talca, broker y outbox, guías hacia ERP y SII, identidad y telemetría crítica. Evita adquirir satélites y negociar el túnel durante la falla, sin costo adicional por el plan de tarifa plana. En los cross-docking es el camino principal, con LTE de dos proveedores como respaldo (ADR-02, Anexo 4-O). \\
  B-02 --- Gateway IoT & Gateway industrial Moxa UC-8200 o equivalente con Greengrass V2, doble entrada DC de 12 a 48 V y dos fuentes de 24 VDC & Talca 2 y Concepción 1; reserva en Talca & 3 + 1 de reserva = 4 & Lee las cámaras por Modbus y conserva 24 h de datos; en Talca una fuente viene del tablero respaldado por UPS y generador, y en Concepción la alimentación se respalda con la UPS del gabinete. \\
  \multicolumn{5}{@{}l}{\intersemibold Infraestructura de la sala técnica de Talca} \\
UPS & Modular de 15 kVA, doble conversión en línea, en N+1 y con bypass de mantenimiento & Zona de UPS y baterías (plano del recinto, apartado 4.3.1.4) & 1 & Dimensionado sobre la carga TI de diseño de 8,9 kW, que incluye el servidor del ERP, y $\approx$ 9,4 kVA aparentes (factor de potencia 0,95), al 80 \% de utilización, $\approx$ 11,8 kVA (RT-06.11; Bases Técnicas Transversales, cap. 6, p. 15); autonomía de 30 min o más a plena carga (RT-06.07; Bases Técnicas Transversales, cap. 6, p. 15). \\
  Generador & Grupo electrógeno de 25 kVA con estanque para 24 h continuas y reabastecimiento contratado & Exterior, junto a la acometida & 1 & Sostiene una carga total de diseño de $\approx$ 13,5 kW ($\approx$ 16,9 kVA, 68 \% de su capacidad) con margen de operación conforme a RT-06.08 (Bases Técnicas Transversales, cap. 6, p. 15). \\
  Transferencia automática & Transferencia red-generador con tablero de protecciones & Acometida, accesible desde el exterior & 1 & Encadena empalme, tablero, transferencia, UPS, PDU y equipo sin interrupción (RT-06.07; Bases Técnicas Transversales, cap. 6, p. 15). La transferencia y el generador se prueban mensualmente con carga real; las instalaciones eléctricas se revisan y miden semestralmente, con informe al CLIENTE (RT-06.10; Bases Técnicas Transversales, cap. 6, p. 15). \\
  Climatización de precisión & 2 unidades de 10 kW ($\approx$ 34.000 BTU/h) en N+1 & Sala de servidores y comunicaciones & 2 & Cubren la carga térmica de diseño de 8,9 kW aun con la pérdida de una unidad (RT-06.13, RT-06.11; Bases Técnicas Transversales, cap. 6, p. 15). \\
  Climatización de la sala de UPS y baterías & Sistema propio de 3,5 kW con ventilación y control de 20--25\,\textdegree{}C & Sala de UPS y baterías de Talca & 1 & Disipa las pérdidas de la UPS y el calor de carga de baterías en un recinto separado de la sala de servidores. \\
  Distribución eléctrica & 2 PDU verticales A/B por gabinete, con medidor & Racks R01 y R02 & 4 & Dos caminos eléctricos por gabinete (RT-08.04; Bases Técnicas Transversales, cap. 8, p. 18); la medición a la salida de las PDU aporta el denominador de la PUE declarada (RT-06.11; Bases Técnicas Transversales, cap. 6, p. 15). \\
  Detección de incendio & Detección temprana por aspiración con tecnología láser, tipo AnaLASER o equivalente & Sala de servidores y comunicaciones & 1 & Detecta antes de que exista llama visible (RT-06.16; Bases Técnicas Transversales, cap. 6, p. 15); junto con la extinción se integra al monitoreo en línea y notifica al NOC y al CLIENTE (RT-06.19; Bases Técnicas Transversales, cap. 6, p. 16). \\
  Monitoreo ambiental del recinto & Sensores en red de temperatura y humedad, y cables detectores de agua bajo el piso técnico y junto a la climatización, con alertas a CloudWatch & Sala de servidores y comunicaciones, y sala de UPS y baterías & 5 sensores de temperatura y humedad (2 por rack y 1 en la sala de UPS) y 2 detectores de agua & Monitorea en línea temperatura, humedad y presencia de agua, con alertamiento integrado a la plataforma de observabilidad (RT-06.14; Bases Técnicas Transversales, cap. 6, p. 15). \\
  Extinción & Agente limpio tipo FM-200 o equivalente, con aprobación UL, instalación conforme a NFPA, botón de aborto y extintores portátiles habilitados & Sala de servidores y comunicaciones & 1 & Extinción automática sobre equipos energizados (RT-06.17; Bases Técnicas Transversales, cap. 6, p. 15) y extintores con mantención y certificación vigentes (RT-06.18; Bases Técnicas Transversales, cap. 6, p. 15). \\
  Control de acceso & Biometría facial con AFIS como respaldo, esclusa antipassback con nueva verificación, bitácora electrónica y estación de enrolamiento interna y externa & Acceso del recinto técnico & 1 & RT-06.20 (Bases Técnicas Transversales, cap. 6, p. 16) a RT-06.23 (Bases Técnicas Transversales, cap. 6, p. 16); bitácora auditable conservada 5 años o más (RT-06.21, Bases Técnicas Transversales, cap. 6, p. 16; RT-16.10, Bases Técnicas Transversales, cap. 16, p. 29). \\
  Videovigilancia & Cámaras IP con imágenes en línea de 30 días o más y grabaciones anteriores en respaldo recuperable y auditable & Acceso principal, pasillo de control, esclusa, sala de servidores y comunicaciones, zona de UPS y baterías, y exterior & 6 & Cubre cada zona del plano del recinto y su perímetro, integrada al monitoreo (RT-06.24; Bases Técnicas Transversales, cap. 6, p. 16). \\
  Recinto de custodia & Recinto con iluminación sin UV de hasta 300 lux, 40 a 60 \% HR, ventilación forzada y 18 a 27\,\textdegree{}C & Área de respaldo del recinto (plano del recinto, apartado 4.3.1.4) & 1 & Conserva los medios de respaldo sin degradarlos (RT-06.27; Bases Técnicas Transversales, cap. 6, p. 16). \\
  Gabinetes & R01 servidores (3 nodos 2U en U18--U23, NAS 2U en U1--U2, KVM 1U en U24, servidor del ERP 2U en U25--U26, conmutador de transferencia 1U en U27, paneles OM4/Cat6A en U40--U42) y R02 comunicaciones (bandeja del operador con ONT de fibra y router LTE, 2U en U31--U32; 2 firewalls 1U en U33--U34; switch de gestión 1U en U35; organizadores 1U en U36 y U39; 2 conmutadores de núcleo 1U en U37--U38; ODF y paneles en U40--U42), 42U & Sala de servidores y comunicaciones & 2 & Racks de servidores independientes de los racks de comunicaciones (RT-06.05; Bases Técnicas Transversales, cap. 6, p. 15); ocupación proyectada de R01 de 15U/42U (36 \%) y de R02 de 12U/42U (29 \%), con 27U y 30U libres para el margen de crecimiento a tres años, coherente con la reserva eléctrica del 20 \% de la tabla de carga eléctrica del apartado 4.3.1. \\
  Piso técnico y cableado & Piso técnico y recorridos de cableado Cat6A y fibra OM4 hacia los paneles y terminaciones & Sala & 1 & Los recorridos se verifican para 10 GbE y cada enlace se documenta y certifica (RT-06.04; Bases Técnicas Transversales, cap. 6, p. 14). \\
  \multicolumn{5}{@{}l}{\intersemibold Gabinetes de borde} \\
UPS de gabinetes de piso & UPS en línea de 1,5 kVA por gabinete, 30 min a plena carga & Talca 4 y Concepción 2 & 6 & Cada UPS alimenta el switch de acceso PoE y la impresora de andén cercana durante cortes breves; la carga de diseño de 0,88 kVA usa el 59~\% de la UPS (Anexo 4-W). \\
Gabinete de borde de Concepción & Gabinete industrializado con UPS en línea de 5 kVA y 30 min a plena carga, sin generador & CD Concepción & 1 & Protege los dos servidores, los firewalls, switches de núcleo, Starlink y gateway IoT, con una carga de diseño de 3,42 kVA que usa el 68~\% de la UPS (Anexo 4-W); sin red eléctrica la bodega no opera por falta de iluminación y equipos de carga, mientras la UPS cubre cortes breves y apagado ordenado y los datos ya se replican fuera del sitio (Bases Técnicas Transversales, cap. 6, numeral 6.1, p. 14). \\
  Climatización del borde & Climatización de precisión acorde al equipamiento del borde & CD Concepción, gabinete de borde & 1 & Dimensionada al equipamiento real del borde, conforme a la tipología del numeral 6.1 (Bases Técnicas Transversales, cap. 6, p. 14). \\
  Control de acceso y monitoreo del borde & Acceso controlado y monitoreo remoto integrado al NOC & CD Concepción, gabinete de borde & 1 & Aplica la tipología de borde operacional del numeral 6.1 con supervisión remota del sitio (Bases Técnicas Transversales, cap. 6, p. 14). \\
  Gabinete de borde de cross-docking & Gabinete de 12U con UPS interna de 0,75 kVA y 30 min, segundo circuito directo de red, cerradura biométrica, sensores ambientales y ventilación & Curicó, Chillán y Los Ángeles & 3 & Los dos circuitos alimentan por separado las fuentes redundantes de los dos mini-PC, los firewalls y los switches conforme al numeral 6.1 (Bases Técnicas Transversales, cap. 6, p. 14); la carga de diseño de 0,42 kVA usa el 57~\% de la UPS (Anexo 4-W). \\
  \multicolumn{5}{@{}l}{\intersemibold Plataforma de nube y servicios gestionados} \\
N-01 a N-03 --- Portales & S3 privado, CloudFront, AWS WAF y API REST pública; el Portal de Clientes se publica además como aplicación web progresiva instalable en el teléfono del cliente & sa-east-1, borde público & 1 aplicación para 3 portales & Publica los portales de clientes, transportistas y proveedores sin exponer los sitios del CLIENTE (RT-03.04; Bases Técnicas Transversales, cap. 3, p. 8). El portal instalado es el perfil de autoatención: pedido, estado de entrega, documentos y saldo, con armado del pedido sin señal (RT-16.30 y RT-17.01 del caso; Bases Técnicas del caso, cap. 15, p. 27); su uso es opcional (ADR-20, Anexo 4-O). \\
  API Gateway & Dos API REST, regional pública y privada, endpoint de interfaz execute-api y autorizador Lambda & Producción en sa-east-1; réplica reducida en us-east-1 & 2 API REST y 1 función Lambda por región & WAF, modelos, cuotas, límites de tasa y VPC Link al ALB privado; acceso público solo por CloudFront (ADR-13). \\
  AWS Verified Access & Confianza OIDC Keycloak y postura CrowdStrike Falcon; AWS WAF de instancia & Producción en sa-east-1; réplica reducida en us-east-1 & 1 instancia por región & Controla por aplicación las consolas y la administración de Keycloak con MFA y postura (ADR-16). \\
  N-04 --- Plataforma de aplicación & ECS Fargate, Application Load Balancer, Elastic Container Registry y CodeBuild; las imágenes se replican entre registros hacia us-east-1 & sa-east-1 en 2 zonas; réplica reducida en us-east-1 & API: 2 a 4 tareas, techo 8; consumidor: 2 a 4; trabajos: 2 a 4; planificador: 1 & La API llega a 7 tareas en la sensibilidad de 91,70 solicitudes/s. Ejecuta los perfiles de API, consumidor, trabajos y planificador del monolito Laravel; erp-sync opera en VM-04. La réplica reducida usa imágenes ya replicadas a ECR de us-east-1 (ADR-01, ADR-09 y ADR-12, Anexo 4-O). \\
  N-04 --- Frontend de consolas & Angular en ECS Fargate tras ALB privado & VPC de Producción, en 2 zonas de sa-east-1 & 2 tareas en 2 zonas & Sirve las consolas y reenvía sus llamadas a la API REST privada por execute-api (RT-03.02; Bases Técnicas Transversales, cap. 3, p. 8; ADR-20, Anexo 4-O). \\
  N-04 --- Motor de optimización de rutas & Contenedor propio en ECS Fargate, tras un contrato versionado con M4 & sa-east-1; relanzamiento en otra zona & 1 tarea por corrida, bajo demanda & Calcula la secuencia de rutas de M4 fuera del artefacto PHP. Si falla, ECS la relanza en otra zona y la corrida se repite dentro de la planificación de 15:00 a 18:30; la aceptación comprueba menos de 20 minutos por corrida (ADR-01 y ADR-12, Anexo 4-O). \\
  N-04 --- Transporte AS2 & OpenAS2 BSD en ECS Fargate tras Network Load Balancer con paso TLS directo & sa-east-1 en 2 zonas & 2 tareas en 2 zonas; 1 acuerdo por cadena & TLS 1.3 en el contenedor, IP registradas, firma, cifrado, certificados por cadena y MDN; prueba de interoperabilidad por cadena (RT-03.02; Bases Técnicas Transversales, cap. 3, p. 8; ADR-11, Anexo 4-O). \\
  N-05 --- Base de datos en nube & Aurora PostgreSQL con Aurora Global Database, instancias db.r6g.2xlarge de 8 vCPU y 64 GiB, y AWS DMS & Clúster primario en sa-east-1 y secundario en us-east-1 & 3 instancias: escritor y lector en sa-east-1, y 1 en us-east-1 & Transaccional de los módulos en nube y réplica de lectura del WMS, separada del estado central que actualizan los eventos; conmuta entre zonas en menos de 30 s y mantiene RPO de 15 min para la réplica del WMS por los tres caminos de Talca (ADR-04 y ADR-09, Anexo 4-O). \\
  N-06 --- Telemetría cruda & DynamoDB bajo demanda; Global Tables solo para temperatura, y la tabla de posiciones de flota solo en sa-east-1 & Temperatura en sa-east-1 y us-east-1; posiciones solo en sa-east-1 & 2 tablas: temperatura global en ambas regiones y posiciones regional & Recibe las lecturas de IoT sin gestión de capacidad y las expira a los 30 días (ADR-04). \\
  N-07 --- Caché & ElastiCache for Redis con primario y réplica & sa-east-1 en 2 zonas & 1 clúster & Responde la consulta de stock y crédito de preventa en menos de 2 s. \\
  N-08 --- Ingesta de IoT & IoT Core con su motor de reglas e IoT Greengrass & sa-east-1 y gateways de borde & 3 gateways y 28 termógrafos & Recibe, valida y alerta la telemetría de cadena de frío, y gestiona la flota de gateways. \\
  N-09 --- Mensajería & SQS FIFO de reconciliación (grupos por sitio y agregado), de solicitudes al ERP (grupos por despacho o documento) y de respuesta para Concepción y cada cross-docking, y de coordinación de reservas para cada uno de los cinco sitios; colas SQS de trabajos de Laravel (notificaciones y EDI); y SNS & sa-east-1; equivalentes vacíos creados por infraestructura como código en us-east-1 & Por región: 13 colas (1 de reconciliación, 1 de solicitudes al ERP, 4 de respuesta, 5 de coordinación de reservas y 2 de trabajos), cada una con su cola de fallidos, y 1 tema SNS & Lleva por conexión saliente las solicitudes al ERP hasta erp-sync de VM-04, devuelve las respuestas a cada sitio, entrega a cada shipper las solicitudes de retención y liberación de stock, y reconcilia eventos idempotentes. Al conmutar de región, los shippers y erp-sync se redirigen a las colas de us-east-1 (ADR-05, ADR-09 y ADR-17, Anexo 4-O). \\
  N-10 --- Plataforma analítica & S3, Glue, Redshift Serverless y QuickSight embebido & sa-east-1; copia diferida en us-east-1 & 1 lago de 5 años & Lectura y autoría con permisos separados en consolas; los indicadores actuales usan la API de negocio. \\
  N-11 --- Respaldo inmutable y réplica & AWS Backup con Vault Lock, S3 Object Lock y replicación de S3 entre regiones & sa-east-1 y us-east-1 & 1 política & Mantiene la copia inmutable del esquema 3-2-1-1-0 y la réplica de recuperación (ADR-09 y ADR-18, Anexo 4-O). La frecuencia, la retención y el tiempo de restauración por dominio de datos (RT-07.13; Bases Técnicas Transversales, cap. 7, p. 18) son los de la tabla de respaldo y restauración por dominio de datos del Subdocumento 4 (apartado 4.2.4). \\
  N-12 --- Seguridad y gobierno & KMS, Secrets Manager, Parameter Store, IAM Identity Center, Shield Advanced, GuardDuty Runtime Monitoring para ECS Fargate, Security Hub, Inspector, Macie, Security Lake, CloudTrail, Config, CloudWatch, Systems Manager, Control Tower, Organizations, Cost Explorer y Budgets & sa-east-1 y us-east-1; claves, secretos, telemetría de runtime y hallazgos en las regiones habilitadas & 8 cuentas: 5 ambientes, gestión, archivo de registros y auditoría & GuardDuty administra el agente sidecar de las tareas nuevas de Fargate y analiza eventos de proceso, archivos y red; Security Hub consolida los hallazgos para el flujo de respuesta de seguridad (Amazon Web Services, s.~f.-a, s.~f.-b). \\
  D-01 --- Firewall y Customer Gateway (parte en nube) & VPC, Transit Gateway, Site-to-Site VPN, NAT Gateway, VPC Endpoints y Route 53 & sa-east-1 y us-east-1 & 2 túneles por sitio y región: 10 hacia sa-east-1 y 10 preparados hacia us-east-1 & Es la contraparte de los mismos túneles contados en la fila D-01 de túneles VPN, sin suma adicional. Conecta los sitios con la nube por VPN cifrada (RT-03.21; Bases Técnicas Transversales, cap. 3, p. 9). \\
  N-13 --- Gestión de dispositivos & Android Enterprise y Zebra DNA como servicio gestionado & Servicio en nube y agente en el dispositivo & 433 equipos: 77 EC55, 121 TC58e y 235 MC9400, incluida la reserva y el crecimiento del año 3 & Mantiene inventario central, configura, actualiza firmware y aplicación, bloquea y borra de forma remota el parque de terreno y bodega (RT-03.18; Bases Técnicas Transversales, cap. 3, p. 9); las reservas se entregan precargadas (Bases Técnicas Transversales, cap. 8, p. 19). \\
  F-03 --- Detección y respuesta en endpoints y cargas de trabajo & CrowdStrike Falcon para servidores y endpoints; GuardDuty Runtime Monitoring para tareas Linux Amazon ECS Fargate, plataforma 1.4 o posterior y rol de ejecución habilitado & Falcon en los 11 equipos físicos, 14 VM y 192 estaciones; GuardDuty en tareas Fargate de sa-east-1 y us-east-1 & 25 servidores + 192 estaciones = 217 endpoints, más agente administrado por tarea Fargate & GuardDuty incorpora el sidecar administrado a tareas nuevas cubiertas y envía los hallazgos a Security Hub; las tareas existentes se reinician o redespliegan al habilitarlo (RT-11.16; Amazon Web Services, s.~f.-a, s.~f.-b). El servidor del ERP del CLIENTE no recibe el agente, porque se traslada sin cambios de software; el firewall de Talca solo le permite el tráfico de la ACL en VM-04 y la administración del CLIENTE. \\
  Dependencias de salida & MDM, EDR, SII, Transbank, GIS, notificaciones y telemetría de flota & Conexiones salientes de la solución & 7 dependencias & Integraciones externas sin entrada a la VPC de Producción (superficie de exposición, apartado 4.2.5). \\
  \multicolumn{5}{@{}l}{\intersemibold Software de base y licenciamiento} \\
A-01 --- Motor WMS (artefacto Laravel) & Monolito modular en Laravel 13 sobre PHP 8.5, con PHP-FPM y PHP CLI, dependencias fijadas por \texttt{composer.lock}; código abierto & Nube (Fargate), VM-01, VM-C01, E-01, VM-03, VM-C04 y VM-04 & 1 imagen firmada & El perfil \texttt{wms\_only} es el Motor WMS de VM-01, VM-C01 y E-01. Una sola base de código y esquema para los perfiles de API, \texttt{wms\_only}, shipper, consumidor, trabajos, planificador y \texttt{erp-sync}; VM-04 ejecuta dos procesos PHP CLI (ADR-01 y ADR-17, Anexo 4-O). \\
  Aplicación de terreno & Aplicación nativa Android en Kotlin, con SQLite y Zebra DataWedge & EC55, TC58e y MC9400 & 1 artefacto para 3 perfiles & Controla los periféricos industriales y opera un turno completo sin señal (ADR-07). \\
  Portales web y consolas & Angular con TypeScript y Tailwind CSS, código abierto & Portales en S3 y CloudFront; consolas en Fargate & 3 portales y consolas & Comparte componentes y contratos TypeScript; las consolas se sirven tras ALB privado. \\
  A-02 --- Base transaccional del WMS & PostgreSQL 16 con PostGIS, código abierto & VM-02, VM-C02 y su copia, y los seis E-01 & 9 instancias, 4 en espera & Registro local de cada bodega durante los cortes. En Concepción y en los cross-docking, la instancia activa se replica en forma sincrónica a la de espera. Solo VM-02 (Talca) se replica a Aurora por DMS; Concepción y los cross-docking sincronizan eventos por los flujos de integración de 4.2.5. \\
  A-03 --- Broker de colas & RabbitMQ, código abierto, con el adaptador AMQP \texttt{php-amqplib} & VM-03, VM-C04 y su copia, y los seis E-01 & 9 instancias, 4 en espera & Buffer de 24 h por sitio con mensajes persistentes y confirmación de publicación; el shipper vacía el buffer hacia SQS FIFO como sobres JSON (ADR-05). \\
  A-04 --- Capa anticorrupción del ERP & Adaptador del artefacto Laravel con dos procesos PHP CLI de \texttt{erp-sync}; código propio, sin licencia de terceros & VM-04, CD Talca & 1 instancia con 2 procesos & Traduce los contratos de la solución al ERP de 2017 sin modificarlo; \texttt{erp-sync} lee la cola local de Talca y, por salida, la cola FIFO de solicitudes al ERP, con una clave idempotente por operación (ADR-17, Anexo 4-O). \\
  A-05 --- Identidad & Keycloak, código abierto, con verificador local de relevo de turno & Fargate, VM-05, VM-C03 y su copia, y los seis E-01 & 1 maestro, 9 cachés y 9 verificadores & Autoridad única en nube con cachés de solo lectura; el verificador local valida el manifiesto de turno firmado y el segundo factor local, el PIN personal, en los relevos sin enlace (ADR-06). \\
  Virtualización & Proxmox VE con Ceph de tres réplicas, código abierto & Clúster de Talca y servidores de Concepción & 1 clúster de 3 nodos y 2 servidores & Talca tolera la pérdida de un nodo y conserva quórum con \texttt{min\_size=2}; los dos servidores de Concepción no forman clúster: cada uno usa RAID 10 local, y la redundancia entre ambos es la réplica sincrónica de PostgreSQL. \\
  Contenedores en los sitios & Docker y Docker Compose, código abierto & Máquinas virtuales de Talca y Concepción y los tres cross-docking & 5 sitios & Ejecutan la misma imagen de la aplicación que la nube; en los cross-docking orquestan además la base, el broker y la caché de identidad, con instalación y parcheo sin dependencia de la red. \\
  F-01 --- Telemetría de observabilidad & AWS Distro for OpenTelemetry y agente SSM & VM-06, VM-C04 y su copia, y los seis E-01 & 9 colectores & Reciben las trazas y métricas de OpenTelemetry para PHP y las guardan hasta 24 h antes de enviarlas a CloudWatch al reconectar (ADR-14). \\
  F-02 --- Gestión de parches & Ansible con líneas base CIS y Terraform para la infraestructura de nube, código abierto & Repositorio del CLIENTE; nodos de los cinco sitios por la red de gestión & 1 repositorio & Aplica parches y líneas base CIS a los nodos de los cinco sitios sin licencia de plataforma, y declara infraestructura y configuración como código versionado (ADR-21, Anexo 4-O). \\
  Integración y entrega continuas & GitLab CI como orquestador y AWS CodeBuild para la construcción & Pipeline del repositorio del CLIENTE; construcción en sa-east-1 & 1 pipeline & Instala las dependencias desde \texttt{composer.lock}, compila, ejecuta PHPUnit, PHPStan/Larastan, Pint y las pruebas de contrato OpenAPI y AsyncAPI, analiza composición con \texttt{composer audit} y escanea secretos e imágenes con bloqueo automático del despliegue (RT-04.05; Bases Técnicas Transversales, cap. 4, p. 10). Construye y firma cada imagen de forma hermética, con procedencia SLSA nivel 3. \\
\end{tablalafrox}

La tabla vincula el equipamiento de sitio con los servicios de nube y los controles de acceso del catálogo de emplazamiento.

\section*{Referencias}

\begin{list}{}{\setlength{\leftmargin}{1.27cm}\setlength{\itemindent}{-1.27cm}\setlength{\itemsep}{0.5\baselineskip}\raggedright}
\item Amazon Web Services. (s.~f.-a). \textit{GuardDuty Runtime Monitoring}. \url{https://docs.aws.amazon.com/guardduty/latest/ug/runtime-monitoring.html}
\item Amazon Web Services. (s.~f.-b). \textit{How Runtime Monitoring works with Fargate (Amazon ECS only)}. \url{https://docs.aws.amazon.com/guardduty/latest/ug/how-runtime-monitoring-works-ecs-fargate.html}
\item Distribuidora Puelche S.A. (2026a). \textit{Bases Administrativas de Licitación N.º TFEP-01/2026}.
\item Distribuidora Puelche S.A. (2026b). \textit{Bases Técnicas Transversales de Licitación N.º TFEP-01/2026}.
\item Distribuidora Puelche S.A. (2026c). \textit{Caso 02: Logística. Especificaciones del problema y operación de Distribuidora Puelche S.A.} (Licitación N.º TFEP-01/2026).
\end{list}


\end{document}
